\section*{Bab \ref{ch:b1:electrophilic-additions} --- Alkena dan Alkuna: Adisi Elektrofilik}

\begin{solution}{exo:b1:electrophilic-additions:1}
2-Kloropropana; 2-bromo-2-metilpropana; 1-bromo-1-metilsikloheksana;
2-klorobut-2-ena.
\end{solution}

\begin{solution}{exo:b1:electrophilic-additions:2}
1,2-Dibromoetana; \textit{trans}-1,2-dibromosikloheksana;
1,2-dibromopropana.
\end{solution}

\begin{solution}{exo:b1:electrophilic-additions:3}
Propan-2-ol; 2-metilpropan-2-ol; 1-metilsikloheksan-1-ol (Markovnikov
dalam setiap kasus).
\end{solution}

\begin{solution}{exo:b1:electrophilic-additions:4}
Propuna (\omterm{def:b1:electrophilic-additions:acetylide}{alkuna terminal}): \ce{CH3C#CH + NH2- -> CH3C#C- + NH3}. Etanol:
\ce{CH3CH2OH + NH2- -> CH3CH2O- + NH3}. But-2-una dan propena tidak
mempunyai hidrogen yang cukup asam.
\end{solution}

\begin{solution}{exo:b1:electrophilic-additions:5}
Pasangan $\pi$ mengambil \ce{H+} pada karbon \ce{CH2}: kation tersier
\ce{(CH3)3C+}; \omterm{def:b1:lewis-resonance-vsepr:lewis}{sepasang elektron bebas} air berikatan dengannya:
\ce{(CH3)3COH2+}; sebuah molekul air mengambil sebuah proton:
\ce{(CH3)3COH}. Dalam asam pekat yang panas, air sedikit dan alkena yang
mudah menguap lepas: reaksi baliknya, \omterm{def:b1:alcohol-activation:dehydration}{dehidrasi}, berlangsung.
\end{solution}

\begin{solution}{exo:b1:electrophilic-additions:6}
Protonasi 2-metilpropena menghasilkan kation tersier, protonasi propena
kation sekunder. Protonasi yang menjadi penentu laju mempunyai \omterm{def:b1:elementary-steps:energy-profile}{keadaan
transisi} yang menyerupai kationnya (Hammond): kation tersier yang lebih
stabil dicapai melalui sawar yang lebih rendah.
\end{solution}

\begin{solution}{exo:b1:electrophilic-additions:7}
Ion bromonium menjembatani satu sisi cincin; bromida membukanya dari sisi
lain: kedua atom brom berakhir \textit{trans}. Dibromida \textit{trans}
itu \omterm{def:b1:stereochemistry-in-depth:stereocentre}{kiral} (tidak ada bidang cermin), tetapi bromida menyerang kedua karbon
sama banyak: kedua \omterm{def:b1:stereochemistry-in-depth:enantiomers}{enantiomer} terbentuk dalam jumlah yang sama, dan
produknya rasemat, tidak aktif optis.
\end{solution}

\begin{solution}{exo:b1:electrophilic-additions:8}
Protonasi C1 menghasilkan kation sekunder pada C2, di sebelah C3 tersier
yang membawa sebuah hidrogen; hidrogen itu berpindah bersama pasangannya
ke C2 (pergeseran hidrida-1,2), sehingga tersisa kation tersier pada C3,
yang ditangkap oleh \ce{Cl-}: 2-kloro-2-metilbutana.
\end{solution}

\begin{solution}{exo:b1:electrophilic-additions:9}
Propuna + natrium amida: ion propunida; dengan 1-bromopropana (SN2):
\ce{CH3C#C-CH2CH2CH3}, heks-2-una. Propunida dengan benzaldehida, lalu
asam encer: \ce{C6H5CH(OH)C#CCH3}, 1-fenilbut-2-un-1-ol.
\end{solution}

\begin{solution}{exo:b1:electrophilic-additions:10}
\cip{E}-but-2-ena: bromonium pada satu sisi, bromida dari sisi lain pada
C2 atau pada C3; kedua serangan menghasilkan senyawa yang sama, dengan
deskriptor \cip{2R,3S}: \omterm{def:b1:stereochemistry-in-depth:enantiomers}{senyawa meso}, satu \omterm{def:b1:stereochemistry-in-depth:isomers}{stereoisomer}. \cip{Z}-but-2-ena:
serangan pada C2 menghasilkan \cip{2S,3S} (atau bayangan cerminnya dari
sisi lain), serangan pada C3 \cip{2R,3R}; jumlahnya sama: dua
\omterm{def:b1:stereochemistry-in-depth:isomers}{stereoisomer}, sebuah \omterm{def:b1:stereochemistry-in-depth:enantiomers}{campuran rasemat}.
\end{solution}

\begin{solution}{exo:b1:electrophilic-additions:11}
Dalam ion bromonium yang tidak simetris, karbon yang lebih tersubstitusi
membawa lebih banyak muatan positif (ikatan C--Br-nya lebih panjang dan
lebih lemah, seperti kation yang hampir tersier atau sekunder). Air, yang
sangat berlebih, menyerang karbon itu dari sisi yang berseberangan dengan
jembatan: 1-bromopropan-2-ol.
\end{solution}

\begin{solution}{exo:b1:electrophilic-additions:12}
$D = 1$: sebuah alkena. 2-Metilbut-1-ena dan 2-metilbut-2-ena sama-sama
menghasilkan kation tersier, sehingga terbentuk 2-bromo-2-metilbutana dan
2-metilbutan-2-ol. Spektrum NMR proton keduanya berbeda: dua proton vinil
(\ce{=CH2}, singlet) untuk yang pertama, satu (\ce{=CH-}, sebuah kuartet)
untuk yang kedua.
\end{solution}

\begin{solution}{pb:b1:electrophilic-additions:1}
\textbf{1.} Pada setiap karbon \ce{CH3} berperingkat lebih tinggi daripada
H: \cip{Z} mempunyai kedua gugus metil pada sisi yang sama, \cip{E} pada
sisi yang berlawanan.
\textbf{2.} \omterm{def:b1:stereochemistry-in-depth:isomers}{Stereoisomer} yang bukan bayangan cermin: \omterm{def:b1:stereochemistry-in-depth:enantiomers}{diastereomer}.
\textbf{3.} Rotasi di sekitar C=C akan memutus ikatan $\pi$, yang
memerlukan energi jauh lebih besar daripada yang diberikan tumbukan pada
suhu kamar.
\textbf{4.} \cip{E}: gugus-gugus metilnya tidak saling berdesakan.
\textbf{5.} \ce{C4H8 + Br2 -> C4H8Br2}.
\textbf{6.} Dua \omterm{def:b1:stereochemistry-in-depth:stereocentre}{pusat stereogenik} yang ekuivalen: paling banyak empat,
sebenarnya tiga (\cip{2R,3R}, \cip{2S,3S}, dan meso \cip{2R,3S}).
\textbf{7.} Pasangan $\pi$ menyerang satu atom Br dari \ce{Br2}, pasangan
Br--Br pergi sebagai \ce{Br-}, dan \omterm{def:b1:lewis-resonance-vsepr:lewis}{sepasang elektron bebas} atom brom yang
diserang berikatan dengan karbon kedua: ion bromonium beranggota tiga.
\textbf{8.} Dalam ion berjembatan, setiap atom memiliki oktet; \omterm{def:b1:electronic-effects:intermediates}{karbokation}
sekunder terbuka akan mempunyai karbon dengan enam elektron.
\textbf{9.} \omterm{def:b1:lewis-resonance-vsepr:lewis}{Sepasang elektron bebas} \ce{Br-} berikatan dengan salah satu
karbon cincin; ikatan C--Br jembatan putus, dan pasangannya berpindah ke
atom brom penjembatan.
\textbf{10.} Jembatan itu menempati satu sisi; seperti pada SN2,
\omterm{def:b1:electronic-effects:nucleophile}{nukleofil} masuk berseberangan dengan ikatan yang putus.
\textbf{11.} Sebuah \omterm{def:b1:electrophilic-additions:halonium}{adisi anti}.
\textbf{12.} Air akan bersaing dengan bromida untuk ion bromonium dan
menghasilkan bromohidrin.
\textbf{13.} \cip{2R,3S} (atau \cip{2S,3R}, senyawa yang sama).
\textbf{14.} Senyawa yang sama.
\textbf{15.} Dalam \omterm{def:b1:stereochemistry-in-depth:isomers}{konformasi} dengan kedua atom brom eklips, sebuah
bidang cermin melewati bagian di antara C2 dan C3; deskriptornya
berlawanan.
\textbf{16.} \cip{2R,3R} dan \cip{2S,3S}.
\textbf{17.} Jumlahnya sama: \omterm{def:b1:stereochemistry-in-depth:enantiomers}{campuran rasemat}.
\textbf{18.} Ya: kedua alkena yang \omterm{def:b1:stereochemistry-in-depth:enantiomers}{diastereomer} itu menghasilkan produk
yang berbeda (meso lawan rasemat).
\textbf{19.} Tidak ada: kedua \omterm{def:b1:stereochemistry-in-depth:enantiomers}{enantiomer} saling meniadakan.
\textbf{20.} $5.60/56.0 = \qty{0.100}{mol}$ masing-masing.
\textbf{21.} $0.100 \times 0.85 \times 215.8 = \qty{18.3}{g}$ masing-masing.
\textbf{22.} \textbf{$[\alpha] = 0^\circ$ (\omterm{def:b1:stereochemistry-in-depth:enantiomers}{senyawa meso}), \qty{18.3}{g}}.
\end{solution}
